% ----------------------------------------------------------------------
%  3.1.5 (продолжение). Металлическая связь и структуры металлов
% ----------------------------------------------------------------------
\subsubsection{Металлическая связь}

В ковалентной связи обобществляются два электрона двух соседних атомов.
В металле обобществлены \emph{все} валентные электроны всех атомов: они не
привязаны к отдельным атомам, а принадлежат кристаллу целиком и
свободно перемещаются по всему объёму. Это \term{свободные электроны},
или \term{электронный газ}.

\begin{opredelenie}
\term{Металлическая связь}~— связь между положительными ионными остовами
атомов и отрицательным электронным газом из коллективизированных
валентных электронов, который как бы «цементирует» ионный остов
(рис.~\ref{fig:metal-bond}).
\end{opredelenie}

\begin{figure}[ht]
\centering
\begin{tikzpicture}[scale=0.6]
  \fill[fmBlue!10] (-0.7,-0.7) rectangle (6.7,4.7);
  \foreach \x in {0,...,6}{\foreach \y in {0,...,4}{
    \node[circle, fill=fmRed!35, draw=fmRed!70!black, minimum size=5.6mm, inner sep=0, font=\scriptsize] at (\x,\y) {$+$};}}
  \foreach \x/\y in {0.5/0.4,1.6/0.6,2.4/1.5,3.5/0.45,4.6/1.45,5.5/0.6,0.4/2.5,1.5/3.4,2.6/2.6,3.4/3.6,4.5/2.4,5.6/3.5,0.6/1.55,6.4/2.5,3.5/1.55,2.5/4.4,4.4/4.45,1.45/2.45,5.45/1.55}{
    \fill[fmBlue] (\x,\y) circle (0.11);}
\end{tikzpicture}
\caption{Модель металла: положительные ионные остовы в узлах решётки и
электронный газ (синие точки), свободно движущийся между ними}
\label{fig:metal-bond}
\end{figure}

Существование свободных электронов доказывает опыт Мандельштама и
Папалекси (и Толмена--Стюарта): при резкой остановке быстро вращающейся
катушки из металлической проволоки в ней возникает кратковременный ток~—
электроны по инерции продолжают двигаться. Об этом же говорят высокие
электро- и теплопроводность металлов.

\paragraph{Почему металлы~— плотные упаковки}
Отдав в общее пользование 1--3 валентных электрона, атом металла
превращается в ион с заполненными оболочками $s$ и $p$, т.\,е. со
сферически симметричным облаком (заполнение внутренних $d$- и
$f$-оболочек картину не меняет). Поэтому
\begin{glavnoe}
Металлическая связь \vazhno{ненаправленна} и \vazhno{ненасыщаема}.
Энергия минимальна, когда атомы упакованы как можно плотнее: атом
металла можно считать шаром, а кристаллическую структуру металла~—
плотнейшей упаковкой шаров.
\end{glavnoe}

\paragraph{Плотные упаковки}
В плоскости шары плотнейшим образом укладываются единственным способом:
каждый касается шести соседей. В пространстве плотнейших упаковок
бесконечно много, но важнейшие две (п.~1.4 и п.~2.2.6):
\begin{itemize}
\item \emph{двухслойная} $ABAB\ldots$~— \term{ГПУ}, отношение $c/a=1{,}633$,
базис $(000;\ \frac13\frac23\frac12)$;
\item \emph{трёхслойная} $ABCABC\ldots$~— \term{ГЦК}, базис
$(000;\ 0\frac12\frac12;\ \frac120\frac12;\ \frac12\frac120)$.
\end{itemize}
Кроме них, металлам свойственна \term{ОЦК}-структура (базис
$(000;\ \frac12\frac12\frac12)$). Она лишь немного уступает плотнейшим
упаковкам, хотя её координационное число 8, а не 12 (рис.~\ref{fig:metal-cells},
табл.~\ref{tab:packing}).

\begin{figure}[ht]
\centering
\begin{subfigure}[b]{0.3\textwidth}\centering
\begin{tikzpicture}[cubeproj, scale=2.1]
  \CubeEdges \FCCback \FCCmid \FCCfront
\end{tikzpicture}
\caption{ГЦК}
\end{subfigure}\hfill
\begin{subfigure}[b]{0.3\textwidth}\centering
\begin{tikzpicture}[cubeproj, scale=2.1]
  \CubeEdges
  \foreach \p in {(0,1,0),(1,1,0),(0,1,1),(1,1,1)} \node[atom=3.4mm] at \p {};
  \node[atomB=3.4mm] at (0.5,0.5,0.5) {};
  \foreach \p in {(0,0,0),(1,0,0),(0,0,1),(1,0,1)} \node[atom=3.4mm] at \p {};
\end{tikzpicture}
\caption{ОЦК}
\end{subfigure}\hfill
\begin{subfigure}[b]{0.34\textwidth}\centering
\begin{tikzpicture}[x={(1cm,0cm)}, y={(0cm,0.4cm)}, z={(0cm,1cm)}, scale=1.05, line join=round]
  \def\cc{1.633}
  % рёбра призмы: задние — пунктиром
  \draw[atomLine, line width=0.5pt, dashed] (1,0,0) -- (0.5,0.866,0) -- (-0.5,0.866,0) -- (-1,0,0);
  \draw[atomLine, line width=0.5pt, dashed] (0.5,0.866,0) -- (0.5,0.866,\cc) (-0.5,0.866,0) -- (-0.5,0.866,\cc);
  \draw[atomLine, line width=0.5pt] (-1,0,0) -- (-0.5,-0.866,0) -- (0.5,-0.866,0) -- (1,0,0);
  \draw[atomLine, line width=0.5pt] (1,0,\cc) -- (0.5,0.866,\cc) -- (-0.5,0.866,\cc) -- (-1,0,\cc)
     -- (-0.5,-0.866,\cc) -- (0.5,-0.866,\cc) -- cycle;
  \foreach \a in {0,180,240,300}{\draw[atomLine, line width=0.5pt] ({cos(\a)},{sin(\a)},0) -- ({cos(\a)},{sin(\a)},\cc);}
  % атомы — от дальних к ближним
  \foreach \zz in {0,\cc}{\foreach \a in {60,120}{\node[atom=3mm] at ({cos(\a)},{sin(\a)},\zz) {};}}
  \node[atomB=3mm] at (0,0.577,\cc/2) {};
  \foreach \zz in {0,\cc}{\node[atom=3mm] at (0,0,\zz) {};
    \foreach \a in {0,180}{\node[atom=3mm] at ({cos(\a)},{sin(\a)},\zz) {};}}
  \foreach \a in {210,330}{\node[atomB=3mm] at ({0.577*cos(\a)},{0.577*sin(\a)},\cc/2) {};}
  \foreach \zz in {0,\cc}{\foreach \a in {240,300}{\node[atom=3mm] at ({cos(\a)},{sin(\a)},\zz) {};}}
  \draw[dimline] (1.35,0,0) -- node[right, font=\small]{$c$} (1.35,0,\cc);
  \draw[dimline] (-0.5,-0.866,-0.35) -- node[below, font=\small]{$a$} (0.5,-0.866,-0.35);
\end{tikzpicture}
\caption{ГПУ}
\end{subfigure}
\caption{Три основные структуры металлов. В ОЦК оранжевым выделен атом в
центре куба; в ГПУ~— атомы среднего слоя (положения $B$ между слоями $A$)}
\label{fig:metal-cells}
\end{figure}

\begin{table}[ht]
\centering
\caption{Плотность заполнения пространства шарами}
\label{tab:packing}
\small
\begin{tabular}{@{}lcc@{}}
\toprule
Решётка & Координационное число & Коэффициент заполнения \\
\midrule
ГПУ & 12 & 0,74 \\
ГЦК & 12 & 0,74 \\
ОЦК & 8 & 0,68 \\
Простая кубическая & 6 & 0,52 \\
Алмаза & 4 & 0,34 \\
Теллура & 2 & 0,23 \\
\bottomrule
\end{tabular}
\end{table}

\begin{primer}[Как считают коэффициент заполнения]
Возьмём ребро куба равным 1. Коэффициент заполнения~— суммарный объём
шаров $\frac43\pi r^3$ в ячейке.
\begin{itemize}
\item \emph{ГЦК}: атомы касаются вдоль диагонали грани ($\sqrt2=4r$),
$r=\sqrt2/4$; в ячейке $8\cdot\frac18+6\cdot\frac12=4$ атома;
$4\cdot\frac43\pi\bigl(\frac{\sqrt2}{4}\bigr)^3=0{,}74$.
\item \emph{ОЦК}: касание вдоль диагонали куба ($\sqrt3=4r$), $r=\sqrt3/4$;
в ячейке $8\cdot\frac18+1=2$ атома; заполнение $0{,}68$.
\item \emph{Простая кубическая}: $r=1/2$, один атом: $0{,}52$.
\item \emph{Алмаз}: касание вдоль четверти диагонали куба, $r=\sqrt3/8$;
8 атомов: $0{,}34$.
\end{itemize}
\end{primer}

Почему ОЦК лишь немного «рыхлее» плотнейших упаковок? У ОЦК вторая
координационная сфера (6 атомов на расстоянии $a$) всего на 15\,\%
дальше первой (8 атомов на $a\sqrt3/2$), а у ГЦК~— на 41\,\% ($a$ против
$a/\sqrt2$). Поэтому иногда условно считают, что у ОЦК-решётки
координационное число 14 ($8+6$).

\subsubsection{Кристаллические структуры металлов по периодической системе}

Более 3/4 элементов~— металлы, и более 3/4 из сотни с лишним
металлических модификаций (включая полиморфные) кристаллизуются в
решётках ГЦК, ГПУ и ОЦК, встречающихся примерно одинаково часто.
Остальные имеют менее симметричные решётки примерно 12 типов.
Периодический закон проявляется и здесь.

\paragraph{Непереходные металлы}
\begin{itemize}
\item I группа: все щелочные металлы~— ОЦК; Cu, Ag, Au~— ГЦК.
\item II группа: Be, Mg, Zn, Cd~— ГПУ, причём реальные решётки могут
заметно отклоняться от идеальной. У Mg $c/a=1{,}624$ (почти идеально
$1{,}633$), у Be $1{,}568$ (сжата вдоль оси $c$), у Zn и Cd~— $1{,}856$ и
$1{,}886$ (вытянуты). У ртути~— ромбоэдрическая решётка с к.\,ч.~6.
\item Ca и Sr~— ГЦК, Ba и Ra~— ОЦК, как щелочные металлы (их внешние
$s$-электроны связаны так же слабо: первые потенциалы ионизации Ba и Ra
меньше, чем у Li).
\item III и следующие группы: закономерности нет. Ga~— сложная
ромбическая, In~— искажённая ГЦК ($c/a=1{,}08$), Tl~— ГПУ. Олово: ниже
$13{,}2$\,°C~— серое (решётка алмаза), выше~— белое (тетрагональная,
искажённая решётка алмаза, $c/a=0{,}55$, к.\,ч.~6). Pb~— типичная ГЦК.
$\alpha$-Po~— простая кубическая (заполнение 0,52).
\end{itemize}

\paragraph{Переходные металлы}
Переходные металлы почти без исключений кристаллизуются в ГЦК, ГПУ и ОЦК
(рис.~\ref{fig:tm-structures}). Плотнейшие упаковки встречаются у
начальных и конечных элементов каждого ряда, ОЦК~— в середине (группы V
и VI). Области существования структур отвечают определённому числу
$[ns+(n-1)d]$-электронов на атом.

\begin{figure}[ht]
\centering
\begin{tikzpicture}[x=12.6mm, y=-9mm]
  \tikzset{cell/.style={draw=white, line width=1.2pt, minimum width=12.4mm, minimum height=8.6mm, inner sep=0, align=center, font=\small}}
  \colorlet{cHCP}{fmBlue!35}\colorlet{cBCC}{fmAmber!45}\colorlet{cFCC}{fmGreen!40}\colorlet{cCX}{fmGray!35}
  \foreach \g [count=\i] in {III,IV,V,VI,VII,VIII,VIII,VIII,I\,B}{\node[font=\scriptsize] at (\i,0.3) {\g};}
  \foreach \el/\st [count=\i] in {Sc/HCP,Ti/HCP,V/BCC,Cr/BCC,Mn/CX,Fe/BCC,Co/HCP,Ni/FCC,Cu/FCC}{
    \node[cell, fill=c\st] at (\i,1) {\textbf{\el}};}
  \foreach \el/\st [count=\i] in {Y/HCP,Zr/HCP,Nb/BCC,Mo/BCC,Tc/HCP,Ru/HCP,Rh/FCC,Pd/FCC,Ag/FCC}{
    \node[cell, fill=c\st] at (\i,2) {\textbf{\el}};}
  \foreach \el/\st [count=\i] in {La/HCP,Hf/HCP,Ta/BCC,W/BCC,Re/HCP,Os/HCP,Ir/FCC,Pt/FCC,Au/FCC}{
    \node[cell, fill=c\st] at (\i,3) {\textbf{\el}};}
  \node[font=\scriptsize, anchor=east] at (0.45,1) {$3d$}; \node[font=\scriptsize, anchor=east] at (0.45,2) {$4d$}; \node[font=\scriptsize, anchor=east] at (0.45,3) {$5d$};
  \begin{scope}[shift={(1,3.9)}]
    \foreach \st/\lab [count=\k from 0] in {HCP/ГПУ, BCC/ОЦК, FCC/ГЦК, CX/сложная}{
      \node[cell, fill=c\st, minimum width=8mm, minimum height=4.5mm] at (\k*2.1,0) {};
      \node[font=\small, anchor=west] at (\k*2.1+0.4,0) {\lab};}
  \end{scope}
\end{tikzpicture}
\caption{Структуры переходных металлов при комнатной температуре
(La~— двойная гексагональная упаковка $ABAC$, отнесена к гексагональным)}
\label{fig:tm-structures}
\end{figure}

При полиморфизме у предплавильной (высокотемпературной) модификации
обычно ОЦК-решётка. Так, у реакторных металлов Ti, Zr, Hf
низкотемпературная модификация~— ГПУ, высокотемпературная~— ОЦК.
Исключение~— железо: у него ОЦК-модификация существует и при низких
температурах ($\alpha$-Fe), и перед плавлением ($\delta$-Fe), а между ними~—
ГЦК ($\gamma$-Fe).

Марганец~— характерное исключение: у него четыре модификации:
\begin{enumerate}[label=\arabic*)]
\item $\alpha$-Mn (ниже $\sim700$\,°C)~— сложная кубическая решётка с 58
атомами в ячейке;
\item $\beta$-Mn ($700$--$1079$\,°C)~— сложная кубическая, 20 атомов в ячейке;
\item $\gamma$-Mn ($1079$--$1143$\,°C)~— гранецентрированная тетрагональная;
\item $\delta$-Mn ($1143$\,°C~— $\Tpl=1244$\,°C)~— ОЦК.
\end{enumerate}

\paragraph{Лантаноиды}
Лантаноиды кристаллизуются в основном в плотные упаковки, но с большей
\emph{слойностью}: в них сочетаются фрагменты ГЦК и ГПУ чередования. Так,
$\alpha$-La, $\beta$-Ce, $\alpha$-Pr, $\alpha$-Nd имеют четырёхслойную
упаковку $ABAC\,ABAC\ldots$ (половина слоёв «кубические», половина
«гексагональные»; по сути, дефект упаковки в каждом четвёртом слое),
$\alpha$-Sm~— девятислойную $ABABCBCAC\ldots$ Из 13 лантаноидов 9
полиморфны, и у всех предплавильная модификация~— ОЦК.

\paragraph{Актиноиды: уран и плутоний}
У актиноидов уровни $5f$, $6d$ и $7s$ близки по энергии; при изменении
температуры электроны перераспределяются между ними, и решётка
перестраивается. Поэтому структуры сложные, а число модификаций от тория
к плутонию растёт (табл.~\ref{tab:U-Pu}).

\begin{table}[ht]
\centering
\caption{Модификации важнейших реакторных металлов~— урана и плутония}
\label{tab:U-Pu}
\small
\begin{tabular}{@{}llll@{}}
\toprule
Фаза & $T$, °C & Решётка & Атомов в ячейке \\
\midrule
$\alpha$-U & до 668 & ромбическая (гофрированные слои) & 4 \\
$\beta$-U & 668--774 & сложная тетрагональная & 30 \\
$\gamma$-U & 774--1132 & ОЦК & 2 \\
\midrule
$\alpha$-Pu & до 120 & моноклинная & 16 \\
$\beta$-Pu & 120--206 & моноклинная & 34 \\
$\gamma$-Pu & 206--319 & ромбическая гранецентр. & 8 \\
$\delta$-Pu & 319--452 & ГЦК & 4 \\
$\eta$-Pu ($\delta'$) & 452--476 & тетрагональная объёмноцентр. & 2 \\
$\varepsilon$-Pu & 476--640 & ОЦК & 2 \\
\bottomrule
\end{tabular}
\end{table}

\begin{vnimanie}[Необычный плутоний]
При переходе $\alpha\to\beta$ объём плутония увеличивается на 8,9\,\%~— для
металлов это необычно много. А у $\delta$-Pu коэффициент линейного
расширения \emph{отрицателен} ($\alpha\approx-8{,}6\cdot10^{-6}$~К$^{-1}$):
при нагреве объём уменьшается. Такие скачки объёма создают огромные
технологические трудности при изготовлении и эксплуатации плутониевого
топлива.
\end{vnimanie}

И у U, и у Pu предплавильная модификация~— ОЦК.

\subsubsection{Ковалентный и металлический радиусы}

\paragraph{Ковалентный радиус}
Рентгеноструктурный анализ даёт расстояние между ядрами атомов,
связанных ковалентной связью. Принимают, что у каждого элемента
ковалентный радиус одинаков во всех его ковалентных соединениях:
$r_A+r_B=d_{AB}$. Например, расстояние C--C в алмазе $0{,}154$~нм, значит,
$r_\text{C}=0{,}077$~нм. Длина связи C--Cl в хлористом метиле
$0{,}176$~нм даёт $r_\text{Cl}=0{,}176-0{,}077=0{,}099$~нм~— и это хорошо
совпадает с половиной расстояния Cl--Cl в молекуле Cl$_2$. Такие расчёты
приближённы: связь никогда не бывает чисто ковалентной.

Радиус сильно зависит от типа связи: в кристаллах As, Sb, Bi расстояние
между атомами, связанными ковалентно, заметно меньше, чем между атомами,
связанными силами Ван-дер-Ваальса (табл.~\ref{tab:AsSbBi}).

\paragraph{Металлический радиус}
\term{Металлический радиус}~— половина кратчайшего расстояния между атомами
металла. Он почти всегда \emph{больше} ковалентного (табл.~\ref{tab:radii}):
металлическая связь заметно отличается от ковалентной.

\begin{table}[ht]
\centering
\caption{Металлические (приведённые к к.\,ч.~12) и ковалентные радиусы, нм}
\label{tab:radii}
\small
\begin{tabular}{@{}lcccccc@{}}
\toprule
 & Li & Be & Na & Ca & Sc & Fe \\
\midrule
$r_\text{мет}$ & 0,155 & 0,112 & 0,190 & 0,197 & 0,162 & 0,126 \\
$r_\text{ков}$ & 0,134 & 0,090 & 0,154 & 0,174 & 0,144 & 0,120 \\
\bottomrule
\end{tabular}
\end{table}

\paragraph{Радиус зависит от координационного числа}
Изменяется ли диаметр атома при полиморфном превращении? Пусть он
постоянен. Тогда при переходе ГЦК (заполнение 0,74) $\to$ ОЦК (0,68) объём
на атом вырос бы в $0{,}74/0{,}68=1{,}09$ раза~— на 9\,\%. При превращении
$\gamma$-Fe (ГЦК) $\to$ $\alpha$-Fe (ОЦК) это, вероятно, приводило бы к
растрескиванию. На опыте объём меняется лишь примерно на 1\,\%: почти
всё уменьшение компактности компенсируется уменьшением межатомных
расстояний.

\begin{glavnoe}
С уменьшением координационного числа межатомные расстояния (атомные
диаметры) уменьшаются, с увеличением~— растут. Поэтому в таблицах
атомные диаметры приводят к одному координационному числу~— 12 с помощью
поправочных коэффициентов (табл.~\ref{tab:goldschmidt}).
\end{glavnoe}

\begin{table}[ht]
\centering
\caption{Относительный атомный диаметр при разных координационных числах
(тип связи не меняется)}
\label{tab:goldschmidt}
\small
\begin{tabular}{@{}lcccccc@{}}
\toprule
Координационное число & 12 & 8 & 6 & 4 & 3 & 1 \\
Атомный диаметр & 1,00 & 0,97 & 0,96 & 0,88 & 0,81 & 0,72 \\
\bottomrule
\end{tabular}
\end{table}

Если же при превращении меняется \emph{тип связи}, объём меняется сильно.
При переходе белого олова (металлическая связь, к.\,ч.~6) в серое
(ковалентная, к.\,ч.~4) атомный радиус уменьшается на 8\,\%, но объём на
атом растёт более чем на 20\,\%. В новых кристаллах возникают огромные
напряжения, хрупкое серое олово рассыпается в порошок~— это
\term{«оловянная чума»}.

То же в твёрдых растворах: хром (ОЦК, к.\,ч.~8), растворяясь в никеле (ГЦК,
к.\,ч.~12), увеличивает свой диаметр, что искажает решётку никеля и
увеличивает её параметр.

\subsubsection{Полиморфизм}

\begin{opredelenie}
\term{Полиморфизм}~— изменение типа кристаллической структуры при
изменении температуры или давления. Разные структуры одного вещества~—
\term{полиморфные модификации}, переход между ними~— \term{полиморфное
превращение}. Оно сопровождается выделением или поглощением скрытой
теплоты, т.\,е. является фазовым переходом первого рода.
\end{opredelenie}

Примеры:
\begin{itemize}
\item цирконий: $\alpha$-Zr (ГПУ, $c/a=1{,}589$) до 863\,°C, $\beta$-Zr (ОЦК)
от 863\,°C до плавления (1855\,°C);
\item щелочные металлы: долго была известна только ОЦК-решётка, но при
очень низких температурах натрий (35~К) и литий (78~К) переходят в
плотноупакованные структуры;
\item под давлением: кремний и германий (полупроводники) переходят в
металлическую модификацию со структурой белого олова.
\end{itemize}

\paragraph{От ковалентной связи к металлической}
Ковалентные кристаллы, построенные по правилу $8-N$, в обычных условиях
не проводят ток как металлы. Но у некоторых из них (Sb, Bi, Si, Ge, Te)
при нагреве вблизи точки плавления появляется металлическая
проводимость, а при плавлении они превращаются в жидкие металлы~—
плавление сопровождается сменой типа связи. У олова такой переход
происходит в твёрдом состоянии (при $13{,}2$\,°C). Низкотемпературные
$\alpha$- и $\beta$-Mn очень хрупки из-за доли направленной ковалентной
связи; при нагреве направленные связи разрушаются, и марганец переходит в
металлические $\gamma$- и $\delta$-модификации. Подобное отмечается у
низкотемпературных модификаций U, Pu, Np.

\begin{glavnoe}
При повышении температуры наблюдаются только переходы от ковалентных
(или ковалентно-металлических) структур к металлическим, но никогда~—
обратно. Полиморфизм особенно распространён у переходных металлов,
лантаноидов и актиноидов: из-за близости уровней $ns$ и $(n-1)d$ (а у
лантаноидов и актиноидов~— ещё и $f$) электроны при изменении
температуры легко переходят между ними. Меняются концентрация
электронного газа и заряд ионных остовов, а вслед за ними~— и
выгодное расположение атомов.
\end{glavnoe}
